\documentclass[12pt]{article}

\usepackage[utf8]{inputenc}
\usepackage[T1]{fontenc}
\usepackage{amsmath, amssymb, amsthm}
\usepackage{geometry}
\geometry{a4paper, margin=2.5cm}

\title{Teorema lui Morera}
\author{}
\date{}

\theoremstyle{plain}
\newtheorem{theorem}{Teoremă}

\begin{document}

\maketitle

\begin{theorem}
\textbf{Teorema lui Morera.}  
Fie $\Omega$ un domeniu deschis din $\mathbb{C}$ și fie $f : \Omega \to \mathbb{C}$ o funcție continuă.  
Presupunem că pentru orice triunghi $T$ complet conținut în $\Omega$ are loc



\[
\int_{\partial T} f(z)\, dz = 0.
\]



Atunci funcția $f$ este olomorfă în $\Omega$.
\end{theorem}

\end{document}
